\section{Sémantika Datalogu}
Zoberme si veľmi jednoduchú teóriu: 
\begin{verbatim}
    p = {(1)}.    
\end{verbatim}

Pre túto teóriu existuje niekoľko modelov (dokonca ich je nekonečne veľa). Najjednoduchší z nich je
$\inter \models_\true \verb|p|(1)$ a pre všetky ostatné formuly je interpretácia nepravdivá. Ďalšími
modelmi môžu byť napríklad:
\begin{itemize}
\item $\inter' \models_\true \verb|p|(i), \forall i \in \mathbb{N}$ a pre všetky ostatné formuly je
interpretácia nepravdivá,
\item $\inter'' \models_\true \verb|p|(i), \forall i \in \mathbb{R}$ a pre všetky ostatné formuly je
interpretácia nepravdivá,
\item $\inter''' \models_\true \verb|p|(1), \inter''' \models_\true \verb|p|(\text{Pes}), \inter'''
\models_\true p(47), \inter''' \models_\true \verb|p|(\text{Štvorec})$ a pre všetky ostatné formuly
je interpretácia nepravdivá.
\end{itemize}

Ako sémantiku našich Datalogových programov budeme uvažovať \textit{well-founded model}.
Well-founded model je minimálny stabilný trojhodnotový model.

\vline

\textit{Trojhodnotový} model znamená, že interpretácia priraďuje atomickým formulám hodnoty
\textit{true} $\true$, \textit{false} $\false$ a \textit{unknown} $\unknown$.

\vline

\textit{Minimálny} model znamená, že spomedzi všetkých modelov vyberáme ten, ktorý obsahuje najmenej
pravdivých atómov. Inými slovami, chceme odvodiť len tie fakty, ktoré sú naozaj nutné. Ak existuje
viacero takých modelov, vyberieme si ten s najmenším počtom pravdivých atómov.

Príklad: Pre teóriu
\begin{verbatim}
    p = {(1)}.
\end{verbatim}
máme viacero modelov. Interpretácia, kde platí iba $p(1)$, je modelom. Interpretácia, kde platí $p(1)$,
$p(2)$ a $p(3)$, je tiež modelom. Minimálny model je však ten prvý, pretože obsahuje najmenej
pravdivých atómov.

\vline

\textit{Stabilný} model je taký model, ktorý sa nezmení po aplikácii \textit{Gelfond-Lifschitz
Transformation} (GLT). GLT funguje takto:
\begin{enumerate}
\item Vezmeme aktuálnu interpretáciu.
\item Z programu odstránime všetky pravidlá, ktorých negované podmienky sú v tejto interpretácii
pravdivé.
\item Zo zostávajúcich pravidiel odstránime všetky negácie.
\item Pre takto upravený program (bez negácií) nájdeme minimálny model.
\item Ak sa tento minimálny model rovná našej pôvodnej interpretácii, potom je naša interpretácia
stabilný model.
\end{enumerate}

Zoberme si program
\begin{verbatim}
    a.
    b :- \+ c.
    c :- \+ b.
    d :- a, \+ c.
\end{verbatim}

a kandidátnu interpretáciu $\inter$, pre ktorú platí
$$\inter \models_\true \verb|a|, \inter \models_\true \verb|b|, \inter \models_\false \verb|c|,
\inter \models_\true \verb|d|\text{.}$$

Pri aplikácii GLT najprv odstránime tie pravidlá, ktorých negované podmienky sú podľa $\inter$
pravdivé. Pravidlo \verb|c :- \+ b.| sa preto odstráni, keďže \verb|b| je v interpretácii $\inter$
pravdivé. V ostatných pravidlách potom odstránime negácie, a dostaneme program
\begin{verbatim}
    a.
    b.
    d :- a.
\end{verbatim}

Minimálny model takto upraveného programu obsahuje práve atómy \verb|a|, \verb|b| a \verb|d|. Keďže
sa zhoduje s interpretáciou, z ktorej sme vychádzali, daná interpretácia je stabilný model pôvodného
programu.

\vline

\textit{Well-founded} model môžeme chápať ako niečo, čo je spoločné pre všetky stabilné modely.
Funguje nasledovne:
\begin{itemize}
\item atómy, ktoré sú $\true$ vo \textbf{všetkých} stabilných modeloch, označíme ako $\true$,
\item atómy, ktoré sú $\false$ vo \textbf{všetkých} stabilných modeloch, označíme ako $\false$,
\item atómy, ktoré sa v rôznych stabilných modeloch líšia, ponecháme ako $\unknown$.
\end{itemize}

V nasledujúcom príklade ukážeme, ako nájsť well-founded model pomocou iteratívneho procesu. 
Podľa \cite{SemanticsQueries} a \cite{WellFounded} sa postup označuje aj ako iteratívny vzostup.

\vline

\textbf{Algoritmus nájdenia well-founded modelu:}
\begin{enumerate}
\item Začneme interpretáciou, kde sú všetky atómy $\false$.
\item Opakovane aplikujeme GLT transformáciu.
\item Po každej transformácii nájdeme minimálny model upraveného programu.
\item Tento minimálny model použijeme ako novú interpretáciu pre ďalšiu iteráciu.
\item Keď sa interpretácia prestane meniť alebo dostaneme niektorú interpretáciu druhýkrát, iterácia
skončila.
\item Výsledný well-founded model získame z hodnôt, ktoré sa stabilizujú, respektíve oscilujú.
\end{enumerate}

Príklad nájdenia well-founded modelu pre program:
\begin{verbatim}
    a.
    b :- a.
    p :- \+ q.
    q :- \+ p.
    u :- v.
\end{verbatim}

Ako prvú interpretáciu zvolíme $\inter_0$, kde všetky atómy sú nastavené na $\false$. Po aplikácii
GLT dostaneme program
\begin{verbatim}
    a.
    b :- a.
    p.
    q.
    u :- v.
\end{verbatim}
Minimálny model tohto programu je $\inter_1$, kde \verb|a|, \verb|b|, \verb|p| a \verb|q| sú
$\true$, a \verb|u| a \verb|v| sú $\false$.

Po druhej aplikácii GLT dostaneme program
\begin{verbatim}
    a.
    b :- a.
    u :- v.
\end{verbatim}
Minimálny model tohto programu je $\inter_2$, kde \verb|a| a \verb|b| sú $\true$, a \verb|p|,
\verb|q|, \verb|u| a \verb|v| sú $\false$.

Po tretej aplikácii GLT a následnom nájdení minimálneho modelu dostaneme $\inter_3$, ktorá sa
zhoduje s $\inter_1$, čiže iterácia sa zastavila. Vo výslednom well-founded modeli $\inter$ budú
stabilizované atómy \verb|a| a \verb|b| nastavené na $\true$, atómy \verb|u| a \verb|v| nastavené na
$\false$ a atómy \verb|p| a \verb|q| zostanú $\unknown$, keďže oscilujú medzi $\true$ a $\false$.
Celý proces je zhrnutý v nasledujúcej tabuľke:

\begin{center}
\begin{tabular}{c|c|c|c|c|c}
Atóm & $\inter_0$ & $\inter_1$ & $\inter_2$ & $\inter_3$ & $\inter$ \\
\hline
\verb|a| & $\false$ & $\true$ & $\true$ & $\true$ & $\true$ \\
\verb|b| & $\false$ & $\true$ & $\true$ & $\true$ & $\true$ \\
\verb|u| & $\false$ & $\false$ & $\false$ & $\false$ & $\false$ \\
\verb|v| & $\false$ & $\false$ & $\false$ & $\false$ & $\false$ \\
\verb|p| & $\false$ & $\true$ & $\false$ & $\true$ & $\unknown$ \\
\verb|q| & $\false$ & $\true$ & $\false$ & $\true$ & $\unknown$ \\
\end{tabular}
\end{center}

